% GENERATED FILE. Edit canonical lessons, then run npm run generate:books.
% canonical-source-hash: fnv1a64:229ff333cddb7e2c
% canonical-lessons: MR-C352-ullekh-karne, MR-C352-namud-karne, MR-C352-bhar-dene, MR-C352-visesatah, MR-C352-pramukh

\chapter{Mentioning and Stressing}
\label{ch:mr-mentioning-and-stressing}
\hlchaptermodality{352}

\begin{chapteropening}
\textbf{By the end of this chapter:} \emph{I can mention something and stress what matters most.}

Complete the last lesson of chapter 352: I can mention something and stress what matters most.
\end{chapteropening}

\section[\texorpdfstring{\mr{उल्लेख} \mr{करणे} (ullekh karṇe)}{उल्लेख करणे (ullekh karṇe)}]{\mr{उल्लेख} \mr{करणे} (ullekh karṇe) — to mention}

\label{lesson:MR-C352-ullekh-karne}



\begin{quote}
Before the new one: say the Marathi for nevertheless, then the Marathi for for example.
\end{quote}

\subsection*{You'll want to know: \mr{उल्लेख} \mr{करणे}}
\textbf{\mr{उल्लेख} \mr{करणे}} — \emph{ullekh karṇe} — \textquotedblleft{}to mention\textquotedblright{}.

\begin{quote}
\mr{अरुणने} \mr{पत्रात} \mr{तुमचा} \mr{उल्लेख} \mr{केला}.
\end{quote}

\begin{quote}
\emph{aruṇne patrāt tumcā ullekh kelā.}
\end{quote}

\begin{quote}
Arun mentioned you in his letter.
\end{quote}

Past, with -\mr{ने} (he …): \mr{त्याने} \mr{उल्लेख} \mr{केला}, \emph{tyāne ullekh kelā}.

What is mentioned takes -\mr{चा}: \mr{तुमचा} \mr{उल्लेख}.

\begin{cousinweb}
\textbf{\mr{उल्लेख} \mr{करणे}} (\emph{ullekh karṇe}) joins the tatsama \mr{उल्लेख}, Sanskrit \emph{ullekha}, \textquotedblleft{}causing to appear clearly\textquotedblright{} (Monier-Williams), to \mr{करणे}, \textquotedblleft{}to do\textquotedblright{}.
\end{cousinweb}

\begin{grammarlens}[title={What you've built}]
One more word for giving reasons and opinions.
\end{grammarlens}

\subsection*{Your turn}
\begin{itemize}
\raggedright
  \item \emph{Say it:} \emph{ullekh karṇe}
  \item \emph{Say it:} \emph{ullekh karṇe}, once more
  \item \emph{Say it:} \emph{aruṇne patrāt tumcā ullekh kelā.} — Arun mentioned you in his letter.
  \item \emph{Recall:} say the Marathi for nevertheless, then the Marathi for for example, then say \emph{ullekh karṇe} again
\end{itemize}

\begin{tcolorbox}[breakable,colback=teal!4,colframe=teal!35!black,title={Before you move on}]
What does \mr{उल्लेख} \mr{करणे} mean? (\textquotedblleft{}To mention\textquotedblright{}.) Say it once more.
\end{tcolorbox}
\section[\texorpdfstring{\mr{नमूद} \mr{करणे} (namūd karṇe)}{नमूद करणे (namūd karṇe)}]{\mr{नमूद} \mr{करणे} (namūd karṇe) — to state (formally)}

\label{lesson:MR-C352-namud-karne}



\begin{quote}
Before the new one: say the Marathi for for example, then the Marathi for to mention.
\end{quote}

\subsection*{You'll want to know: \mr{नमूद} \mr{करणे}}
\textbf{\mr{नमूद} \mr{करणे}} — \emph{namūd karṇe} — \textquotedblleft{}to state (formally)\textquotedblright{}.

\begin{quote}
\mr{अहवालात} \mr{सगळे} \mr{खर्च} \mr{नमूद} \mr{केले} \mr{आहेत}.
\end{quote}

\begin{quote}
\emph{ahvālāt sagḷe kharca namūd kele āhet.}
\end{quote}

\begin{quote}
All the expenses are stated in the report.
\end{quote}

Past, with -\mr{ने} (he …): \mr{त्याने} \mr{नमूद} \mr{केलं}, \emph{tyāne namūd kelaṁ}.

Formal, of writing: \mr{अर्जात} \mr{नमूद} \mr{करणे}, to state in an application.

\begin{cousinweb}
\textbf{\mr{नमूद} \mr{करणे}} (\emph{namūd karṇe}) joins \mr{नमूद}, a loan usually traced to Persian \emph{namūd}, \textquotedblleft{}appearance, showing\textquotedblright{}, to \mr{करणे}, \textquotedblleft{}to make\textquotedblright{}.
\end{cousinweb}

\begin{grammarlens}[title={What you've built}]
One more word for giving reasons and opinions.
\end{grammarlens}

\subsection*{Your turn}
\begin{itemize}
\raggedright
  \item \emph{Say it:} \emph{namūd karṇe}
  \item \emph{Say it:} \emph{namūd karṇe}, once more
  \item \emph{Say it:} \emph{ahvālāt sagḷe kharca namūd kele āhet.} — All the expenses are stated in the report.
  \item \emph{Recall:} say the Marathi for for example, then the Marathi for to mention, then say \emph{namūd karṇe} again
\end{itemize}

\begin{tcolorbox}[breakable,colback=teal!4,colframe=teal!35!black,title={Before you move on}]
What does \mr{नमूद} \mr{करणे} mean? (\textquotedblleft{}To state (formally)\textquotedblright{}.) Say it once more.
\end{tcolorbox}
\section[\texorpdfstring{\mr{भर} \mr{देणे} (bhar deṇe)}{भर देणे (bhar deṇe)}]{\mr{भर} \mr{देणे} (bhar deṇe) — to emphasize; to stress}

\label{lesson:MR-C352-bhar-dene}



\begin{quote}
Before the new one: say the Marathi for to mention, then the Marathi for to state.
\end{quote}

\subsection*{You'll want to know: \mr{भर} \mr{देणे}}
\textbf{\mr{भर} \mr{देणे}} — \emph{bhar deṇe} — \textquotedblleft{}to emphasize; to stress\textquotedblright{}.

\begin{quote}
\mr{शिक्षकांनी} \mr{वाचनावर} \mr{भर} \mr{दिला}.
\end{quote}

\begin{quote}
\emph{śikṣakānnī vācnāvar bhar dilā.}
\end{quote}

\begin{quote}
The teacher stressed reading.
\end{quote}

Past, with -\mr{ने} (he …): \mr{त्याने} \mr{भर} \mr{दिला}, \emph{tyāne bhar dilā}.

The point stressed takes -\mr{वर}: \mr{वाचनावर} \mr{भर} \mr{देणे}.

\begin{cousinweb}
\textbf{\mr{भर} \mr{देणे}} (\emph{bhar deṇe}) joins \mr{भर}, \textquotedblleft{}weight, load\textquotedblright{}, inherited from Sanskrit \emph{bhara}, \textquotedblleft{}carrying, weight\textquotedblright{} (CDIAL 9393), to \mr{देणे}, \textquotedblleft{}to give\textquotedblright{}: to put weight on a point.
\end{cousinweb}

\begin{grammarlens}[title={What you've built}]
One more word for giving reasons and opinions.
\end{grammarlens}

\subsection*{Your turn}
\begin{itemize}
\raggedright
  \item \emph{Say it:} \emph{bhar deṇe}
  \item \emph{Say it:} \emph{bhar deṇe}, once more
  \item \emph{Say it:} \emph{śikṣakānnī vācnāvar bhar dilā.} — The teacher stressed reading.
  \item \emph{Recall:} say the Marathi for to mention, then the Marathi for to state, then say \emph{bhar deṇe} again
\end{itemize}

\begin{tcolorbox}[breakable,colback=teal!4,colframe=teal!35!black,title={Before you move on}]
What does \mr{भर} \mr{देणे} mean? (\textquotedblleft{}To emphasize; to stress\textquotedblright{}.) Say it once more.
\end{tcolorbox}
\section[\texorpdfstring{\mr{विशेषतः} (viśeṣataḥ)}{विशेषतः (viśeṣataḥ)}]{\mr{विशेषतः} (viśeṣataḥ) — especially}

\label{lesson:MR-C352-visesatah}



\begin{quote}
Before the new one: say the Marathi for to state, then the Marathi for to emphasize.
\end{quote}

\subsection*{You'll want to know: \mr{विशेषतः}}
\textbf{\mr{विशेषतः}} — \emph{viśeṣataḥ} — \textquotedblleft{}especially\textquotedblright{}.

\begin{quote}
\mr{हिवाळ्यात}, \mr{विशेषतः} \mr{सकाळी}, \mr{खूप} \mr{थंडी} \mr{असते}.
\end{quote}

\begin{quote}
\emph{hivāḷyāt, viśeṣataḥ sakāḷī, khūp thaṇḍī aste.}
\end{quote}

\begin{quote}
In winter, especially in the mornings, it is very cold.
\end{quote}

Formal; in speech \mr{खास} \mr{करून}.

\begin{cousinweb}
\textbf{\mr{विशेषतः}} (\emph{viśeṣataḥ}) is a tatsama, taken from Sanskrit \emph{viśeṣatas}, \textquotedblleft{}especially, particularly\textquotedblright{} (Monier-Williams).
\end{cousinweb}

\begin{grammarlens}[title={What you've built}]
One more word for giving reasons and opinions.
\end{grammarlens}

\subsection*{Your turn}
\begin{itemize}
\raggedright
  \item \emph{Say it:} \emph{viśeṣataḥ}
  \item \emph{Say it:} \emph{viśeṣataḥ}, once more
  \item \emph{Say it:} \emph{hivāḷyāt, viśeṣataḥ sakāḷī, khūp thaṇḍī aste.} — In winter, especially in the mornings, it is very cold.
  \item \emph{Recall:} say the Marathi for to state, then the Marathi for to emphasize, then say \emph{viśeṣataḥ} again
\end{itemize}

\begin{tcolorbox}[breakable,colback=teal!4,colframe=teal!35!black,title={Before you move on}]
What does \mr{विशेषतः} mean? (\textquotedblleft{}Especially\textquotedblright{}.) Say it once more.
\end{tcolorbox}
\section[\texorpdfstring{\mr{प्रमुख} (pramukh)}{प्रमुख (pramukh)}]{\mr{प्रमुख} (pramukh) — leading; chief}

\label{lesson:MR-C352-pramukh}



\begin{quote}
Before the new one: say the Marathi for to emphasize, then the Marathi for especially.
\end{quote}

\subsection*{You'll want to know: \mr{प्रमुख}}
\textbf{\mr{प्रमुख}} — \emph{pramukh} — \textquotedblleft{}leading; chief\textquotedblright{}.

\begin{quote}
\mr{शहरातील} \mr{प्रमुख} \mr{रस्ते} \mr{आज} \mr{बंद} \mr{आहेत}.
\end{quote}

\begin{quote}
\emph{śahrātīl pramukh raste āj band āhet.}
\end{quote}

\begin{quote}
The main roads in the city are closed today.
\end{quote}

Also a noun, the head: \mr{संस्थेचा} \mr{प्रमुख}.

\begin{cousinweb}
\textbf{\mr{प्रमुख}} (\emph{pramukh}) is a tatsama, taken from Sanskrit \emph{pramukha}, \textquotedblleft{}facing; foremost, chief\textquotedblright{} (Monier-Williams).
\end{cousinweb}

\begin{grammarlens}[title={What you've built}]
One more word for giving reasons and opinions.
\end{grammarlens}

\subsection*{Your turn}
\begin{itemize}
\raggedright
  \item \emph{Say it:} \emph{pramukh}
  \item \emph{Say it:} \emph{pramukh}, once more
  \item \emph{Say it:} \emph{śahrātīl pramukh raste āj band āhet.} — The main roads in the city are closed today.
  \item \emph{Recall:} say the Marathi for to emphasize, then the Marathi for especially, then say \emph{pramukh} again
\end{itemize}

\begin{tcolorbox}[breakable,colback=teal!4,colframe=teal!35!black,title={Before you move on}]
What does \mr{प्रमुख} mean? (\textquotedblleft{}Leading; chief\textquotedblright{}.) Say it once more.
\end{tcolorbox}
